\section{Geometric inputs and Frobenius conventions}
\label{inp:section}

We use the constructible $\ell$-adic geometric Langlands framework.  The
results recalled here apply under Assumption~\ref{intro:characteristic}:
its first three conditions are those of \cite[\S14.4.1]{AGKRRV1}, and its
fourth is the additional condition of \cite[\S D.1.1]{AGKRRV1}.
All categories are $\E$-linear.  A compact object means an object whose
Hom functor commutes with filtered colimits; compactness is stronger than
constructibility on an algebraic stack.

\subsection{Restricted Langlands and functions}

\begin{theorem}[Restricted geometric Langlands]\label{inp:langlands}
There is a Frobenius-invariant union $Z$ of connected components of
$\LocSys^{\mathrm{restr}}_{\Ghat}(\overline X)$ and an equivalence
\[
 \mathbb L_G^{\mathrm{restr}}:\cC
       \xrightarrow{\sim}\IndCoh_{\Nhat}(Z).
\]
It identifies $\cC_Y$ with $\IndCoh_Y(Z)$ for every closed
$\Ghat$-invariant $Y\subset\Nhat$ and intertwines Hecke functors with
tensoring by the corresponding evaluation bundles.  If
$f:Z\to Z$ sends a local system to its pullback along geometric
Frobenius of $\overline X$, the functor corresponding to
$\Phi=(\operatorname{Frob}_{\Bun})_*$ is $f^!$.
The embedding $\cC\hookrightarrow\Shv(\Bun_G(\overline X),\E)$
preserves compact objects.
\end{theorem}

The equivalence and the last assertion are
\cite[Theorems~1.3.9(i) and~1.1.7]{GR}; Frobenius invariance of the
retained union is \cite[\S1.5.4]{GR}.
The Hecke action is compatible with the spectral linear structure
\cite[Lemma~1.3.7]{GR}, and the support compatibility is
\cite[Theorem~4.1.1 and Corollary~4.2.1]{Raskin}.
The stated Frobenius convention is exactly \cite[\S1.3.4]{GLR}.
Thus $Z$ is the prime union of the introduction throughout this paper.
In particular, this theorem does not require identifying that union with
the entire restricted stack.
Under Assumption~\ref{intro:characteristic}, the companion full-support
theorem~\cite[Theorem~1.4, regime~A]{RGLCompanion} identifies
$Z=\LocSys^{\mathrm{restr}}_{\Ghat}(\overline X)$ as spectral prestacks
over the fixed coefficient field $E=\E$.

At a geometric point $x$, a representation $D$ of $\Ghat$ supplies the
vector bundle $\sigma\mapsto D_{\sigma,x}$ on $Z$.  At a closed point of
degree $d$, the geometric points form a Frobenius cycle of length $d$.
The corresponding Hecke construction uses all $d$ evaluation bundles,
with cyclic transport.  This observation will keep the degree of a
closed point visible in our trace calculation.

\begin{remark}[Theta normalization]\label{inp:theta}
The theta normalization of the equivalence can be made over $\mathbf F_q$.
Indeed, $X$ admits a line bundle $\vartheta$ defined over $\mathbf F_q$
with $\vartheta^{\otimes2}\simeq\omega_X$.
For odd characteristic this follows from
\cite[\S5, Remark~2 after Proposition~5.2, p.~61]{Atiyah}: there is a
Frobenius-invariant square-root class, and the Brauer group of the
finite field vanishes, so this class descends.  In characteristic two, choose a rational
function $a\in\mathbf F_q(X)$ with $da\ne0$.  The orders of $da$ are even,
as a Laurent expansion over each perfect residue field shows.  Hence
$\mathcal O_X(\tfrac12\operatorname{div}(da))$ is such a square root
\cite[\S4, p.~191]{Mumford}.  No rational point of $X$ is needed.
\end{remark}

\begin{theorem}[Trace and ordinary local terms]\label{inp:trace}
There is a canonical isomorphism of complexes
\[
 \operatorname{LT}:\Tr(\Phi,\cC)\xrightarrow{\sim}\cA_{\E},
\]
with the right-hand side in degree zero.  For a compact object $M$ and
a morphism $b:M\to\Phi M$, it sends the categorical class $[M,b]$ to
the ordinary alternating trace function of the adjoint morphism
$(\operatorname{Frob}_{\Bun})^*M\to M$.  The function has finite support.
This comparison is compatible with Hecke operations and with the
sheaf--function operations of pullback, star tensor product, and
compactly supported pushforward.
\end{theorem}

The trace isomorphism is \cite[Theorem~0.2.6, equivalently
Corollary~4.1.4]{AGKRRV3}.  Its identification with local terms is
\cite[Theorem~0.6.8, equivalently Theorem~5.2.3, and
Remark~0.6.5]{AGKRRV3}; the compact-embedding hypothesis used there is
provided by Theorem~\ref{inp:langlands}.  The compatibilities and the
stack convention for summation are recalled in \cite[\S1.1]{GLR}.
This is a statement about all of $\cA_{\E}$.  It does not restrict to
cuspidal or Hecke-finite functions.  The stronger assertion that actual
invertible Weil structures span the supported traces will be proved below.

\subsection{Duality and finite evaluations}

\begin{theorem}[Nonstandard duality]\label{inp:duality}
The pairing
\begin{equation}\label{inp:duality-pairing}
 \cC\otimes\cC\longrightarrow\operatorname{Vect}_{\E},\qquad
 (M,N)\longmapsto
 R\Gamma_c\bigl(\Bun_G(\overline X),M\overset{*}{\otimes}N\bigr)
\end{equation}
defines a self-duality of $\cC$.  It is invariant under simultaneous
Frobenius.  If $M,N$ are compact, its value is a perfect complex, also
after any finite Hecke representation is inserted in either variable.
\end{theorem}

The duality is \cite[Theorem~3.2.2]{AGKRRV2}.
Here is the finiteness deduction, which is useful when applying it.
By Theorem~\ref{inp:langlands}, the two objects are compact in ambient
sheaves, hence constructible \cite[Appendix~A.1.6]{AGKRRV2}.
Star tensoring an ambient compact with a constructible object
preserves compactness \cite[Appendix~A.1.6]{AGKRRV2}.
For $p:\Bun_G\to\operatorname{pt}$, the functor $p_!$ is left adjoint
to the continuous functor $p^!$ \cite[Appendix~A.1.7]{AGKRRV2}, so it
preserves compactness.  Its value is therefore compact in
$\operatorname{Vect}_{\E}$, hence perfect.
The Hecke functor for a finite-dimensional representation has a
continuous adjoint supplied by its dual representation, and likewise
preserves compactness.  Frobenius invariance of
\eqref{inp:duality-pairing} follows by change of variables and the
compatibility of Frobenius with star tensor product.

For compact Weil objects the numerical trace of
\eqref{inp:duality-pairing} is consequently
\begin{equation}\label{inp:function-pairing}
 [u,g]=\sum_{b\in S}\frac{u(b)g(b)}{|\Aut(b)|}.
\end{equation}
This is the Grothendieck trace formula with groupoid counting measure;
see \cite[\S\S1.1.9 and~1.2.3]{GLR}.
The pairing is bilinear over $\E$, without complex conjugation.
The compact objects of the categorical dual are the opposites of the
compact objects of $\cC$, and their evaluation is enriched Hom
\cite[\S\S0.5.3 and~5.1.2]{AGKRRV2}.
Thus the second input may be transported to a compact Hom-test object.
Its Weil arrow is reversed on passing to the opposite category.  For
arrows $a:\Phi M\to M$ and $b:N\to\Phi N$, the induced endomorphism of
$\Hom(M,N)$ is
\begin{equation}\label{inp:hom-action}
 h\longmapsto\Phi^{-1}(b\circ h\circ a).
\end{equation}
This convention fixes the direction of every Frobenius action below.
The later spanning theorem will justify testing against every function
in \eqref{inp:function-pairing}.

\subsection{The geometry of a spectral component}

\begin{theorem}[Components and deformation theory]\label{inp:components}
Every connected component $Z_\sigma$ of the restricted local-system
stack has a unique semisimple isomorphism class $\sigma$.  Its
automorphism group $J=\Aut(\sigma)$ is reductive, possibly disconnected.
There is a formal affine coarse space $S_\sigma$ with reduced space a
point, and $Z_\sigma\to S_\sigma$ is relatively algebraic.
After choosing a frame at a geometric point, the component is formal
affine.  For an affine infinitesimal stage $S_n\to S_\sigma$, its
base change is affine; almost finite type stages and bounded Postnikov
truncations may be used.  These affine stages have coherent cohomology
in each bounded range.  The tangent complex at a local system $\tau$ is
\begin{equation}\label{inp:tangent}
 T_\tau\LocSys^{\mathrm{restr}}_{\Ghat}(\overline X)
       \simeq R\Gamma(\overline X,\ghat_\tau)[1].
\end{equation}
In particular the stack is quasi-smooth.
\end{theorem}

The component and finiteness statements assemble
\cite[Theorems~1.4.5 and~1.6.3, Remarks~1.4.6--1.4.7,
Corollary~3.7.4, and Theorem~5.4.2]{AGKRRV1}.
The affine fiber assertion is Corollary~5.4.6 there; the assertion for
general affine infinitesimal stages is made explicitly in
\cite[\S7.9.10]{AGKRRV1}.
The deformation formula is \cite[Proposition~2.2.2 and
\S\S2.4,~25.4.3]{AGKRRV1}.
Formal quasi-smoothness is established in
\cite[\S\S21.2.1--21.2.3]{AGKRRV1}; the deformation complex has the
required amplitude because the curve has cohomological dimension two.
Semisimplicity of $\sigma$ makes its geometric monodromy
reductive, and its centralizer $J$ is reductive in characteristic zero.

For $d$ frames, write $K=\Ghat^d$.  Adding the remaining frames is an
affine torsor, so the same finite-stage assertions hold.  The unique
closed $K$-orbit is $K/J$: closed orbits correspond to semisimple local
systems, and the stabilizer of a framed representative is $J$.
Every nonempty invariant closed subset of a finite stage meets this
orbit.  These statements supply the finite-stage setting for the
reconstruction in the next section; they do not yet identify the
component with a quadratic moment-map model.

\begin{theorem}[Pure Weil lift]\label{inp:pure-lift}
Fix an abstract field isomorphism $\iota:\E\xrightarrow{\sim}\C$.
If $Z_\sigma$ is Frobenius invariant, its semisimple point $\sigma$
admits a Weil structure that is $\iota$-pure of weight zero.  In
particular, writing $\ghat_\sigma$ for its adjoint local system, every
Frobenius eigenvalue on
$H^i(\overline X,\ghat_\sigma)$ has $\iota$-absolute value $q^{i/2}$,
for $0\leq i\leq2$.
\end{theorem}

The pure-lift application to a Frobenius-fixed semisimple point is
stated in \cite[Appendix~A.5.4--A.5.5]{GLR}, using
\cite[Corollary~3.7.4 and Proposition~25.4.6]{AGKRRV1}.
The literal statement of that proposition treats irreducible Weil
parameters; the cited appendix records the semisimple consequence used
here.  Purity of $H^1$ is the smooth proper curve application of
Weil~II \cite[Corollary~3.3.6 and \S3.3.10]{Deligne}, also stated explicitly in
\cite[Appendix~A.5.5]{GLR}.  For $H^0$ it follows by viewing horizontal
sections in a pure fiber, and for $H^2$ by curve duality.
No semisimplicity of the cohomological Frobenius operators is asserted
or required.

Put $\mathfrak j=\Lie(J)$ and
$V=H^1(\overline X,\ghat_\sigma)$.
An invariant form on $\ghat$ and curve duality give
\[
 H^0(\overline X,\ghat_\sigma)=\mathfrak j,
 \qquad H^2(\overline X,\ghat_\sigma)\simeq\mathfrak j^*(-1),
\]
and a $J$-invariant symplectic form
$V\otimes V\to\E(-1)$.  Its moment map is the cup--bracket
obstruction.  After forgetting the Tate factor it is the usual
quadratic map $\mu:V\to\mathfrak j^*$.
We will use purity to linearize Frobenius on the full formal component,
including this obstruction and the evaluation bundles.
